\documentclass[11pt,a4paper]{article}

% --- Encoding and fonts ---
\usepackage[utf8]{inputenc}
\usepackage[T1]{fontenc}

% --- Mathematics ---
\usepackage{amsmath,amssymb,amsthm}
\usepackage{mathtools}

% --- Page layout ---
\usepackage[margin=1in,headheight=14pt]{geometry}

% --- Hyperlinks ---
\usepackage[colorlinks=true,linkcolor=red,citecolor=blue,urlcolor=blue]{hyperref}

% --- Lists ---
\usepackage{enumitem}

% --- Colors and boxes ---
\usepackage{xcolor}
\usepackage[theorems,skins,breakable]{tcolorbox}

% --- Header/footer ---
\usepackage{fancyhdr}
\pagestyle{fancy}
\fancyhf{}
\fancyhead[L]{\textsc{Lesson 9: Convex Optimization and Duality}}
\fancyhead[R]{\thepage}

% --- Tables ---
\usepackage{booktabs}

% ============================================================
% Theorem environments
% ============================================================
\theoremstyle{definition}
\newtheorem{definition}{Definition}[section]
\newtheorem{example}{Example}[section]
\newtheorem{exercise}{Exercise}[section]

\theoremstyle{plain}
\newtheorem{theorem}{Theorem}[section]
\newtheorem{lemma}[theorem]{Lemma}
\newtheorem{proposition}[theorem]{Proposition}
\newtheorem{corollary}[theorem]{Corollary}

\theoremstyle{remark}
\newtheorem{remark}{Remark}[section]

% ============================================================
% Colored boxes
% ============================================================
\newtcolorbox{keyresult}[1][]{
  colback=green!5!white,
  colframe=green!60!black,
  fonttitle=\bfseries,
  title={Key Result},
  breakable,
  #1
}

\newtcolorbox{rlconnection}[1][]{
  colback=blue!5!white,
  colframe=blue!60!black,
  fonttitle=\bfseries,
  title={RL Connection},
  breakable,
  #1
}

\newtcolorbox{warning}[1][]{
  colback=red!5!white,
  colframe=red!60!black,
  fonttitle=\bfseries,
  title={Warning},
  breakable,
  #1
}

% ============================================================
% Custom commands
% ============================================================
\newcommand{\R}{\mathbb{R}}
\newcommand{\N}{\mathbb{N}}
\newcommand{\Q}{\mathbb{Q}}
\newcommand{\Z}{\mathbb{Z}}
\newcommand{\C}{\mathbb{C}}
\newcommand{\Prob}{\mathbb{P}}
\newcommand{\E}{\mathbb{E}}
\newcommand{\indicator}{\mathbf{1}}

\DeclareMathOperator{\dom}{dom}
\DeclareMathOperator{\epi}{epi}
\DeclareMathOperator{\conv}{conv}
\DeclareMathOperator{\prox}{prox}
\DeclareMathOperator*{\argmin}{arg\,min}
\DeclareMathOperator*{\argmax}{arg\,max}
\DeclareMathOperator{\tr}{tr}
\DeclareMathOperator{\diag}{diag}
\DeclareMathOperator{\interior}{int}

% ============================================================
\title{%
  \textbf{Lesson 9: Convex Optimization and Duality} \\[6pt]
  \large Lagrangian Duality, KKT Conditions, and First-Order Methods
}
\author{%
  Session 9 of 102 \quad$\mid$\quad Phase I: Mathematical Foundations \quad$\mid$\quad 8h \\[4pt]
  \textit{Primary texts: Boyd \& Vandenberghe Ch.\,4--5, 9--11; Nesterov Ch.\,1--4}
}
\date{}

\begin{document}
\maketitle
\thispagestyle{fancy}
\tableofcontents
\newpage

%% ============================================================
%% SECTION 1: Introduction
%% ============================================================
\section{Introduction}\label{sec:intro}

Optimization is the computational engine beneath virtually every algorithm in
machine learning and reinforcement learning. Policy gradient methods solve
surrogate optimization problems at each iteration; value function approximation
minimizes projected Bellman error; and constrained MDPs reduce to saddle-point
problems whose structure is governed by duality theory.

This document develops the theory of convex optimization from first principles:
the hierarchy of convex programs (Section~\ref{sec:convex-problems}), Lagrangian
duality and the passage from primal to dual (Section~\ref{sec:duality}), the
Karush--Kuhn--Tucker conditions (Section~\ref{sec:kkt}), gradient descent and
its convergence guarantees (Section~\ref{sec:gd}), Nesterov's accelerated method
(Section~\ref{sec:nesterov}), proximal methods for non-smooth objectives
(Section~\ref{sec:proximal}), and interior-point methods
(Section~\ref{sec:interior}). We close with exercises
(Section~\ref{sec:exercises}) and a summary table of convergence rates
(Section~\ref{sec:summary}).

Throughout, we assume familiarity with real analysis (continuity, compactness)
and linear algebra (eigenvalues, positive definiteness) at the level of the
Phase~00 prerequisites.

%% ============================================================
%% SECTION 2: Convex Optimization Problems
%% ============================================================
\section{Convex Optimization Problems}\label{sec:convex-problems}

\subsection{Standard Form}

\begin{definition}[Convex optimization problem]\label{def:convex-opt}
A \emph{convex optimization problem} in standard form is
\begin{equation}\label{eq:std-form}
\begin{aligned}
  \min_{x \in \R^n} \quad & f_0(x) \\
  \text{subject to} \quad & f_i(x) \le 0, \quad i = 1,\dots,m, \\
                          & a_j^\top x = b_j, \quad j = 1,\dots,p,
\end{aligned}
\end{equation}
where $f_0, f_1, \dots, f_m : \R^n \to \R$ are convex functions and the
equality constraints are affine.
\end{definition}

\begin{definition}[Feasible set, optimal value]\label{def:feasible}
The \emph{feasible set} is
$\mathcal{F} = \{x \in \R^n : f_i(x) \le 0,\; a_j^\top x = b_j\}$.
The \emph{optimal value} is $p^* = \inf\{f_0(x) : x \in \mathcal{F}\}$.
A point $x^*$ with $f_0(x^*) = p^*$ is an \emph{optimal point} (or
\emph{minimizer}).
\end{definition}

\begin{theorem}[Local optimality implies global optimality]\label{thm:local-global}
For the convex program~\eqref{eq:std-form}, every local minimizer is a global
minimizer.
\end{theorem}

\begin{proof}
Let $x^*$ be a local minimizer: there exists $\delta > 0$ such that
$f_0(x^*) \le f_0(x)$ for all feasible $x$ with $\|x - x^*\| \le \delta$.
Suppose for contradiction there exists a feasible $y$ with $f_0(y) < f_0(x^*)$.
Consider $z_t = (1-t)x^* + ty$ for $t \in (0,1)$. By convexity of each $f_i$
and affinity of the equality constraints, $z_t$ is feasible. Moreover,
\[
  f_0(z_t) \le (1-t) f_0(x^*) + t f_0(y) < f_0(x^*)
\]
for all $t \in (0,1)$. Choosing $t$ small enough that
$\|z_t - x^*\| = t\|y - x^*\| < \delta$ contradicts local optimality.
\end{proof}

\subsection{Important Special Cases}

\begin{definition}[Linear program -- LP]\label{def:lp}
A \emph{linear program} has the form
\[
  \min_{x} \; c^\top x \quad \text{s.t.} \quad Ax \le b, \; Cx = d,
\]
where all functions are affine (hence convex).
\end{definition}

\begin{definition}[Quadratic program -- QP]\label{def:qp}
A \emph{quadratic program} minimizes a convex quadratic over a polyhedral
feasible set:
\[
  \min_{x} \; \tfrac{1}{2} x^\top P x + q^\top x \quad \text{s.t.} \quad Ax \le b,
\]
with $P \succeq 0$ (positive semidefinite).
\end{definition}

\begin{definition}[Second-order cone program -- SOCP]\label{def:socp}
An SOCP has constraints of the form $\|A_i x + b_i\|_2 \le c_i^\top x + d_i$
(second-order cone constraints), which are convex.
\end{definition}

\begin{definition}[Semidefinite program -- SDP]\label{def:sdp}
An SDP minimizes a linear objective subject to a linear matrix inequality:
\[
  \min_{x} \; c^\top x \quad \text{s.t.} \quad F_0 + \sum_{i=1}^n x_i F_i \succeq 0,
\]
where $F_0, \dots, F_n$ are symmetric matrices.
\end{definition}

\begin{remark}[Hierarchy]\label{rem:hierarchy}
We have the strict inclusion chain:
$\text{LP} \subset \text{QP} \subset \text{SOCP} \subset \text{SDP} \subset
\text{convex programs}$.
\end{remark}

\begin{rlconnection}
In RL, linear programs arise in the LP formulation of MDPs (minimize
$\sum_s c(s) V(s)$ subject to Bellman feasibility constraints). Semidefinite
programs appear in robust MDP formulations and certain Lyapunov-based safety
constraints. Quadratic programs arise in trust-region policy optimization when
the surrogate objective is approximated to second order.
\end{rlconnection}

%% ============================================================
%% SECTION 3: Lagrangian Duality
%% ============================================================
\section{Lagrangian Duality}\label{sec:duality}

\subsection{The Lagrangian and Dual Function}

\begin{definition}[Lagrangian]\label{def:lagrangian}
Given the convex program~\eqref{eq:std-form}, the \emph{Lagrangian}
$L : \R^n \times \R^m \times \R^p \to \R$ is
\begin{equation}\label{eq:lagrangian}
  L(x, \lambda, \nu) = f_0(x) + \sum_{i=1}^m \lambda_i f_i(x)
    + \sum_{j=1}^p \nu_j (a_j^\top x - b_j),
\end{equation}
where $\lambda \in \R^m_+$ are the \emph{dual variables} (or \emph{Lagrange
multipliers}) for the inequality constraints and $\nu \in \R^p$ are the dual
variables for the equality constraints.
\end{definition}

\begin{definition}[Lagrange dual function]\label{def:dual-fn}
The \emph{Lagrange dual function} $g : \R^m \times \R^p \to
\R \cup \{-\infty\}$ is
\begin{equation}\label{eq:dual-fn}
  g(\lambda, \nu) = \inf_{x \in \R^n} L(x, \lambda, \nu).
\end{equation}
\end{definition}

\begin{proposition}[Dual function is concave]\label{prop:dual-concave}
For any optimization problem (not necessarily convex), the dual function $g$ is
concave, even when the primal problem is not convex.
\end{proposition}

\begin{proof}
For each fixed $x$, the map $(\lambda, \nu) \mapsto L(x, \lambda, \nu)$ is
affine (hence concave) in $(\lambda, \nu)$. The pointwise infimum over $x$ of
a family of concave functions is concave.
\end{proof}

\subsection{Weak Duality}

\begin{definition}[Dual problem]\label{def:dual-problem}
The \emph{Lagrange dual problem} is
\begin{equation}\label{eq:dual-problem}
  \max_{\lambda \ge 0,\, \nu} \; g(\lambda, \nu).
\end{equation}
Its optimal value is denoted $d^*$.
\end{definition}

\begin{theorem}[Weak duality]\label{thm:weak-duality}
For any optimization problem (not necessarily convex),
\begin{equation}\label{eq:weak-duality}
  d^* \le p^*.
\end{equation}
The quantity $p^* - d^*$ is called the \emph{duality gap}.
\end{theorem}

\begin{proof}
Let $\tilde{x}$ be any feasible point for the primal problem, so
$f_i(\tilde{x}) \le 0$ and $a_j^\top \tilde{x} = b_j$. For any
$\lambda \ge 0$ and any $\nu$,
\[
  L(\tilde{x}, \lambda, \nu)
  = f_0(\tilde{x}) + \sum_{i=1}^m \underbrace{\lambda_i}_{\ge 0}
    \underbrace{f_i(\tilde{x})}_{\le 0}
  + \sum_{j=1}^p \nu_j \underbrace{(a_j^\top \tilde{x} - b_j)}_{= 0}
  \le f_0(\tilde{x}).
\]
Taking the infimum over $x$ gives $g(\lambda, \nu) \le L(\tilde{x}, \lambda, \nu) \le f_0(\tilde{x})$.
Since this holds for every feasible $\tilde{x}$, we get $g(\lambda, \nu) \le p^*$.
Maximizing the left side over $(\lambda, \nu)$ with $\lambda \ge 0$ yields
$d^* \le p^*$.
\end{proof}

\subsection{Strong Duality and Slater's Condition}

\begin{definition}[Strong duality]\label{def:strong-duality}
\emph{Strong duality} holds when $d^* = p^*$, i.e., the duality gap is zero.
\end{definition}

\begin{definition}[Slater's condition]\label{def:slater}
A convex program~\eqref{eq:std-form} satisfies \emph{Slater's condition}
(or \emph{Slater's constraint qualification}) if there exists a point
$\hat{x} \in \text{relint}(\dom f_0)$ such that
\[
  f_i(\hat{x}) < 0, \quad i = 1,\dots,m, \qquad
  a_j^\top \hat{x} = b_j, \quad j = 1,\dots,p.
\]
Such $\hat{x}$ is called a \emph{strictly feasible point}.
\end{definition}

\begin{theorem}[Slater implies strong duality]\label{thm:strong-duality}
If the convex program~\eqref{eq:std-form} is feasible and Slater's condition
holds, then strong duality holds: $d^* = p^*$. Moreover, if $p^*$ is finite,
the dual optimum is attained.
\end{theorem}

\begin{proof}[Proof sketch]
We use a supporting-hyperplane argument on the set
\[
  \mathcal{G} = \{(u, v, t) \in \R^m \times \R^p \times \R :
    \exists\, x \text{ with } f_i(x) \le u_i,\; a_j^\top x - b_j = v_j,\;
    f_0(x) \le t\}.
\]
This set is convex (the $f_i$ are convex and $a_j^\top x - b_j$ is affine).
The point $(0, 0, p^*)$ lies on the boundary of $\mathcal{G}$ (it is achieved
in the limit by a minimizing sequence).

We claim $(0, 0, p^*-\epsilon) \notin \mathcal{G}$ for any $\epsilon > 0$, by
definition of $p^*$. Hence there is a supporting hyperplane at $(0, 0, p^*)$:
there exist $(\lambda, \nu, \mu) \ne 0$ with $\mu \ge 0$ such that
\[
  \lambda^\top u + \nu^\top v + \mu t \ge \mu p^*
\]
for all $(u, v, t) \in \mathcal{G}$, with $\lambda \ge 0$.

\emph{Claim:} $\mu > 0$. Suppose $\mu = 0$. Then
$\lambda^\top u + \nu^\top v \ge 0$ for all $(u,v,t) \in \mathcal{G}$.
Evaluating at $x = \hat{x}$ (the Slater point) with $u_i = f_i(\hat{x}) < 0$,
$v_j = 0$, gives $\sum_i \lambda_i f_i(\hat{x}) \ge 0$. Since $\lambda \ge 0$
and $f_i(\hat{x}) < 0$, we get $\lambda = 0$. Then $\nu^\top v \ge 0$ for all
achievable $v$, which forces $\nu = 0$ (since $v$ can be any value via
adjusting $x$). This contradicts $(\lambda, \nu, \mu) \ne 0$.

With $\mu > 0$, divide through by $\mu$ and set $\bar{\lambda} = \lambda/\mu$,
$\bar{\nu} = \nu/\mu$. For any feasible $x$ (where $u_i = f_i(x)$, $v_j = 0$,
$t = f_0(x)$):
\[
  \sum_i \bar{\lambda}_i f_i(x) + f_0(x) \ge p^*.
\]
Taking the infimum over $x$ gives
$g(\bar{\lambda}, \bar{\nu}) \ge p^*$. Combined with weak duality $g \le p^*$,
we conclude $d^* \ge g(\bar{\lambda}, \bar{\nu}) = p^*$, hence $d^* = p^*$.
The dual optimum is attained at $(\bar{\lambda}, \bar{\nu})$.
\end{proof}

\begin{rlconnection}
Strong duality is the theoretical foundation for the Lagrangian approach to
constrained MDPs (CMDPs). In a CMDP, we augment the reward objective with
penalty terms for constraint violations. Slater's condition guarantees there is
no duality gap, so solving the dual (which is unconstrained) is equivalent to
solving the primal. This justifies primal-dual policy optimization algorithms
such as Lagrangian PPO for safe RL.
\end{rlconnection}

%% ============================================================
%% SECTION 4: KKT Conditions
%% ============================================================
\section{Karush--Kuhn--Tucker Conditions}\label{sec:kkt}

\begin{definition}[KKT conditions]\label{def:kkt}
For the convex program~\eqref{eq:std-form}, the \emph{Karush--Kuhn--Tucker
(KKT) conditions} for the primal-dual triple $(x^*, \lambda^*, \nu^*)$ are:
\begin{enumerate}[label=(KKT\arabic*)]
  \item \textbf{Primal feasibility:}
    $f_i(x^*) \le 0$ for $i = 1,\dots,m$; \;
    $a_j^\top x^* = b_j$ for $j = 1,\dots,p$.
  \item \textbf{Dual feasibility:}
    $\lambda^*_i \ge 0$ for $i = 1,\dots,m$.
  \item \textbf{Complementary slackness:}
    $\lambda^*_i f_i(x^*) = 0$ for $i = 1,\dots,m$.
  \item \textbf{Stationarity:}
    $\nabla f_0(x^*) + \sum_{i=1}^m \lambda^*_i \nabla f_i(x^*)
     + \sum_{j=1}^p \nu^*_j a_j = 0$.
\end{enumerate}
\end{definition}

\begin{theorem}[KKT necessary conditions under strong duality]\label{thm:kkt-necessary}
If strong duality holds and $x^*$ is primal optimal, $(\lambda^*, \nu^*)$ is
dual optimal, and the $f_i$ are differentiable at $x^*$, then the KKT
conditions hold at $(x^*, \lambda^*, \nu^*)$.
\end{theorem}

\begin{proof}
Since strong duality holds, $f_0(x^*) = p^* = d^* = g(\lambda^*, \nu^*)$.
By definition,
\[
  g(\lambda^*, \nu^*) = \inf_x L(x, \lambda^*, \nu^*) \le L(x^*, \lambda^*, \nu^*).
\]
We expand:
\[
  L(x^*, \lambda^*, \nu^*)
  = f_0(x^*) + \sum_i \lambda^*_i f_i(x^*) + \sum_j \nu^*_j(a_j^\top x^* - b_j).
\]
By primal feasibility, $a_j^\top x^* - b_j = 0$ and $f_i(x^*) \le 0$. Since
$\lambda^* \ge 0$, we have $\sum_i \lambda^*_i f_i(x^*) \le 0$, so
$L(x^*, \lambda^*, \nu^*) \le f_0(x^*)$.

Combining: $f_0(x^*) = g(\lambda^*, \nu^*) \le L(x^*, \lambda^*, \nu^*) \le f_0(x^*)$.
Hence all inequalities are equalities. In particular,
$\sum_i \lambda^*_i f_i(x^*) = 0$. Since each term is non-positive
($\lambda^*_i \ge 0$, $f_i(x^*) \le 0$), each term is zero:
$\lambda^*_i f_i(x^*) = 0$. This gives \textbf{complementary slackness}.

Furthermore, $g(\lambda^*, \nu^*) = L(x^*, \lambda^*, \nu^*)$, meaning $x^*$
minimizes $L(\cdot, \lambda^*, \nu^*)$ over $\R^n$. Since this is an
unconstrained minimum of a differentiable function, the gradient vanishes:
\[
  \nabla_x L(x^*, \lambda^*, \nu^*)
  = \nabla f_0(x^*) + \sum_i \lambda^*_i \nabla f_i(x^*) + \sum_j \nu^*_j a_j
  = 0.
\]
This is \textbf{stationarity}. Primal and dual feasibility hold by assumption.
\end{proof}

\begin{theorem}[KKT sufficient conditions for convex problems]\label{thm:kkt-sufficient}
For the convex program~\eqref{eq:std-form} with differentiable $f_0, \dots, f_m$,
if a triple $(x^*, \lambda^*, \nu^*)$ satisfies the KKT conditions, then $x^*$
is primal optimal and $(\lambda^*, \nu^*)$ is dual optimal, with zero duality
gap.
\end{theorem}

\begin{proof}
By stationarity, $x^*$ minimizes $L(\cdot, \lambda^*, \nu^*)$ (since
$L$ is convex in $x$ and the gradient vanishes at $x^*$). Therefore,
\[
  g(\lambda^*, \nu^*) = \inf_x L(x, \lambda^*, \nu^*) = L(x^*, \lambda^*, \nu^*).
\]
By complementary slackness and primal feasibility,
\[
  L(x^*, \lambda^*, \nu^*) = f_0(x^*) + \underbrace{\sum_i \lambda^*_i f_i(x^*)}_{=\,0}
  + \underbrace{\sum_j \nu^*_j (a_j^\top x^* - b_j)}_{=\,0} = f_0(x^*).
\]
Hence $g(\lambda^*, \nu^*) = f_0(x^*)$. By weak duality,
$d^* \ge g(\lambda^*, \nu^*) = f_0(x^*) \ge p^* \ge d^*$, so all are equal.
Thus $x^*$ is primal optimal and $(\lambda^*, \nu^*)$ is dual optimal.
\end{proof}

\begin{keyresult}[title={KKT Conditions for Convex Problems}]
For convex programs with differentiable objectives and constraints, the KKT
conditions are \textbf{necessary} (assuming a constraint qualification such as
Slater's condition, which implies strong duality) and \textbf{sufficient} for
optimality. This makes them the central tool for characterizing solutions.
\end{keyresult}

\begin{rlconnection}
The KKT conditions characterize solutions of the policy optimization
subproblem in TRPO and PPO. When we constrain the policy update to a KL
trust region, the KKT multiplier on the KL constraint determines the
effective step size. Complementary slackness tells us: either the constraint
is active (the update reaches the trust-region boundary) or the multiplier is
zero (the unconstrained optimum lies inside the trust region).
\end{rlconnection}

%% ============================================================
%% SECTION 5: Gradient Descent
%% ============================================================
\section{Gradient Descent}\label{sec:gd}

\subsection{Smoothness and Convexity}

\begin{definition}[$L$-smoothness]\label{def:smooth}
A differentiable function $f : \R^n \to \R$ is \emph{$L$-smooth} (or has
$L$-Lipschitz continuous gradient) if
\begin{equation}\label{eq:smooth}
  \|\nabla f(x) - \nabla f(y)\| \le L \|x - y\| \quad \forall\, x, y \in \R^n.
\end{equation}
\end{definition}

\begin{lemma}[Descent lemma]\label{lem:descent}
If $f$ is $L$-smooth, then for all $x, y \in \R^n$:
\begin{equation}\label{eq:descent-lemma}
  f(y) \le f(x) + \nabla f(x)^\top (y - x) + \frac{L}{2}\|y - x\|^2.
\end{equation}
\end{lemma}

\begin{proof}
By the fundamental theorem of calculus,
\[
  f(y) - f(x) = \int_0^1 \nabla f(x + t(y-x))^\top (y-x)\, dt.
\]
Subtracting $\nabla f(x)^\top(y-x)$:
\begin{align*}
  f(y) - f(x) - \nabla f(x)^\top(y-x)
  &= \int_0^1 [\nabla f(x + t(y-x)) - \nabla f(x)]^\top (y-x)\, dt \\
  &\le \int_0^1 \|\nabla f(x + t(y-x)) - \nabla f(x)\| \cdot \|y-x\|\, dt \\
  &\le \int_0^1 L t \|y-x\|^2\, dt = \frac{L}{2}\|y-x\|^2. \qedhere
\end{align*}
\end{proof}

\begin{definition}[$\mu$-strong convexity]\label{def:strong-convex}
A differentiable function $f$ is \emph{$\mu$-strongly convex} (with $\mu > 0$)
if for all $x, y$:
\begin{equation}\label{eq:strong-convex}
  f(y) \ge f(x) + \nabla f(x)^\top(y - x) + \frac{\mu}{2}\|y - x\|^2.
\end{equation}
\end{definition}

\begin{remark}\label{rem:condition-number}
If $f$ is $\mu$-strongly convex and $L$-smooth, the \emph{condition number} is
$\kappa = L/\mu \ge 1$. The condition number governs the convergence speed of
first-order methods.
\end{remark}

\subsection{Gradient Descent Algorithm}

The gradient descent iteration with step size $\eta > 0$ is:
\begin{equation}\label{eq:gd}
  x_{k+1} = x_k - \eta \nabla f(x_k).
\end{equation}

\subsection{Convergence for Convex Functions}

\begin{theorem}[$O(1/k)$ convergence for convex $L$-smooth functions]%
\label{thm:gd-convex}
Let $f$ be convex and $L$-smooth with minimizer $x^*$. Gradient descent with
step size $\eta = 1/L$ satisfies
\begin{equation}\label{eq:gd-convex-rate}
  f(x_k) - f(x^*) \le \frac{L \|x_0 - x^*\|^2}{2k}.
\end{equation}
\end{theorem}

\begin{proof}
Set $\eta = 1/L$. By the descent lemma (Lemma~\ref{lem:descent}) with
$y = x_{k+1} = x_k - \frac{1}{L}\nabla f(x_k)$:
\begin{equation}\label{eq:gd-descent-step}
  f(x_{k+1}) \le f(x_k) + \nabla f(x_k)^\top(x_{k+1} - x_k)
    + \frac{L}{2}\|x_{k+1} - x_k\|^2
  = f(x_k) - \frac{1}{2L}\|\nabla f(x_k)\|^2.
\end{equation}
So the sequence of function values is non-increasing, and
\begin{equation}\label{eq:func-decrease}
  f(x_k) - f(x_{k+1}) \ge \frac{1}{2L}\|\nabla f(x_k)\|^2.
\end{equation}
By convexity, $f(x^*) \ge f(x_k) + \nabla f(x_k)^\top(x^* - x_k)$, hence
\begin{equation}\label{eq:convexity-bound}
  f(x_k) - f(x^*) \le \nabla f(x_k)^\top(x_k - x^*) \le \|\nabla f(x_k)\| \cdot \|x_k - x^*\|.
\end{equation}
Now consider the distance to the optimum:
\begin{align}
  \|x_{k+1} - x^*\|^2
  &= \|x_k - \tfrac{1}{L}\nabla f(x_k) - x^*\|^2 \notag \\
  &= \|x_k - x^*\|^2 - \frac{2}{L}\nabla f(x_k)^\top(x_k - x^*)
     + \frac{1}{L^2}\|\nabla f(x_k)\|^2. \label{eq:dist-recursion}
\end{align}
Using convexity: $\nabla f(x_k)^\top(x_k - x^*) \ge f(x_k) - f(x^*)$, so
\[
  \|x_{k+1} - x^*\|^2 \le \|x_k - x^*\|^2
    - \frac{2}{L}(f(x_k) - f(x^*)) + \frac{1}{L^2}\|\nabla f(x_k)\|^2.
\]
From~\eqref{eq:func-decrease}: $\|\nabla f(x_k)\|^2 \le 2L(f(x_k) - f(x_{k+1}))$, so
\begin{align*}
  \|x_{k+1} - x^*\|^2
  &\le \|x_k - x^*\|^2 - \frac{2}{L}(f(x_k) - f(x^*))
    + \frac{2}{L}(f(x_k) - f(x_{k+1})) \\
  &= \|x_k - x^*\|^2 - \frac{2}{L}(f(x_{k+1}) - f(x^*)).
\end{align*}
Rearranging: $f(x_{k+1}) - f(x^*) \le \frac{L}{2}(\|x_k - x^*\|^2 - \|x_{k+1} - x^*\|^2)$.

Let $\delta_k = f(x_k) - f(x^*)$. Summing from $k = 0$ to $k = K-1$:
\[
  \sum_{k=1}^{K} \delta_k \le \frac{L}{2}\|x_0 - x^*\|^2.
\]
Since $\delta_k$ is non-increasing (from~\eqref{eq:func-decrease}),
$K \delta_K \le \sum_{k=1}^K \delta_k \le \frac{L}{2}\|x_0 - x^*\|^2$, giving
$\delta_K \le \frac{L\|x_0 - x^*\|^2}{2K}$.
\end{proof}

\subsection{Convergence for Strongly Convex Functions}

\begin{theorem}[Linear convergence for strongly convex $L$-smooth functions]%
\label{thm:gd-strong}
Let $f$ be $\mu$-strongly convex and $L$-smooth. Gradient descent with
$\eta = 1/L$ satisfies
\begin{equation}\label{eq:gd-strong-rate}
  f(x_k) - f(x^*) \le \left(1 - \frac{\mu}{L}\right)^k (f(x_0) - f(x^*)).
\end{equation}
Hence $f(x_k) - f(x^*) \le \epsilon$ after $k = O(\kappa \log(1/\epsilon))$
iterations, where $\kappa = L/\mu$.
\end{theorem}

\begin{proof}
From~\eqref{eq:gd-descent-step},
$f(x_{k+1}) \le f(x_k) - \frac{1}{2L}\|\nabla f(x_k)\|^2$.
By strong convexity applied at the optimum $x^*$:
\[
  f(x^*) \ge f(x_k) + \nabla f(x_k)^\top(x^* - x_k) + \frac{\mu}{2}\|x^* - x_k\|^2.
\]
Thus $f(x_k) - f(x^*) \le \nabla f(x_k)^\top(x_k - x^*) - \frac{\mu}{2}\|x_k - x^*\|^2$.
By the Cauchy--Schwarz inequality and the AM-GM inequality (or by optimizing
over $\|x_k - x^*\|$), one obtains the standard inequality for strongly convex
functions:
\begin{equation}\label{eq:PL}
  \|\nabla f(x_k)\|^2 \ge 2\mu (f(x_k) - f(x^*)).
\end{equation}
This is the \emph{Polyak--Lojasiewicz (PL) inequality}, which follows from
strong convexity.

\emph{Proof of~\eqref{eq:PL}:} By strong convexity, for all $y$:
$f(y) \ge f(x_k) + \nabla f(x_k)^\top(y - x_k) + \frac{\mu}{2}\|y - x_k\|^2$.
Minimizing the right side over $y$ (a quadratic in $y$) gives
$y^* = x_k - \frac{1}{\mu}\nabla f(x_k)$, yielding
\[
  f(x^*) \ge f(x_k) - \frac{1}{2\mu}\|\nabla f(x_k)\|^2,
\]
which rearranges to~\eqref{eq:PL}.

Substituting~\eqref{eq:PL} into the descent bound:
\[
  f(x_{k+1}) \le f(x_k) - \frac{1}{2L} \cdot 2\mu(f(x_k) - f(x^*))
  = f(x_k) - \frac{\mu}{L}(f(x_k) - f(x^*)).
\]
Subtracting $f(x^*)$ from both sides:
\[
  f(x_{k+1}) - f(x^*) \le \left(1 - \frac{\mu}{L}\right)(f(x_k) - f(x^*)).
\]
Applying this recursively yields~\eqref{eq:gd-strong-rate} with
$\rho = 1 - \mu/L = 1 - 1/\kappa$.
\end{proof}

\begin{keyresult}[title={Gradient Descent Convergence Rates}]
For $L$-smooth functions with step size $\eta = 1/L$:
\begin{itemize}
  \item \textbf{Convex:} $f(x_k) - f^* = O(1/k)$, i.e., $O(1/\epsilon)$
    iterations for $\epsilon$-accuracy.
  \item \textbf{$\mu$-strongly convex:} $f(x_k) - f^* = O(\rho^k)$ with
    $\rho = 1 - 1/\kappa$, i.e., $O(\kappa \log(1/\epsilon))$ iterations.
\end{itemize}
\end{keyresult}

\begin{rlconnection}
Policy gradient methods and TD learning with function approximation perform
(stochastic) gradient descent on the policy parameters or value function weights.
The convergence rate depends critically on the smoothness ($L$) and strong
convexity ($\mu$) of the loss landscape. In practice, the loss for neural network
policies is non-convex, but locally the landscape often behaves like a strongly
convex function near a good solution.
\end{rlconnection}

%% ============================================================
%% SECTION 6: Nesterov Acceleration
%% ============================================================
\section{Nesterov's Accelerated Gradient Method}\label{sec:nesterov}

\subsection{The Accelerated Iteration}

Nesterov's accelerated gradient method improves the $O(1/k)$ rate for convex
functions to $O(1/k^2)$, which is optimal among first-order methods.

\begin{definition}[Nesterov's accelerated gradient method]\label{def:nesterov}
For a convex, $L$-smooth function $f$, with step size $\eta = 1/L$, the
accelerated gradient method maintains two sequences:
\begin{equation}\label{eq:nesterov}
\begin{aligned}
  y_{k} &= x_k + \frac{k-1}{k+2}(x_k - x_{k-1}), \\
  x_{k+1} &= y_k - \frac{1}{L}\nabla f(y_k),
\end{aligned}
\end{equation}
with $x_0 = x_{-1}$ (so $y_0 = x_0$). The momentum coefficient
$\theta_k = (k-1)/(k+2)$ is the key ingredient.
\end{definition}

\subsection{Convergence Rate}

\begin{theorem}[$O(1/k^2)$ convergence of Nesterov's method]\label{thm:nesterov}
Under the conditions of Definition~\ref{def:nesterov},
\begin{equation}\label{eq:nesterov-rate}
  f(x_k) - f(x^*) \le \frac{2L\|x_0 - x^*\|^2}{(k+1)^2}.
\end{equation}
\end{theorem}

\begin{proof}
We use the \emph{estimate sequence} technique. Define the sequences
$\{A_k\}$ and $\{v_k\}$ by $A_0 = 0$, $a_{k+1} = (k+1)/2$,
$A_{k+1} = A_k + a_{k+1}$, and $v_0 = x_0$. One can verify that
$A_k = k(k+1)/4$ for $k \ge 1$.

We prove by induction the key inequality:
\begin{equation}\label{eq:estimate-seq}
  A_k (f(x_k) - f(x^*)) \le \frac{L}{2}\|v_0 - x^*\|^2
    - \frac{L}{2}\|v_k - x^*\|^2.
\end{equation}
Since $\|v_k - x^*\|^2 \ge 0$, this implies
$f(x_k) - f(x^*) \le \frac{L\|x_0 - x^*\|^2}{2A_k}
= \frac{2L\|x_0 - x^*\|^2}{k(k+1)} \le \frac{2L\|x_0 - x^*\|^2}{(k+1)^2}$.

\emph{Base case} ($k = 0$): $A_0 = 0$, so both sides equal zero.

\emph{Inductive step:} Assume~\eqref{eq:estimate-seq} holds for $k$. We need to
show it for $k+1$. By the descent lemma applied at $y_k$:
\[
  f(x_{k+1}) \le f(y_k) - \frac{1}{2L}\|\nabla f(y_k)\|^2.
\]
By convexity: $f(x^*) \ge f(y_k) + \nabla f(y_k)^\top(x^* - y_k)$, so
$f(y_k) \le f(x^*) + \nabla f(y_k)^\top(y_k - x^*)$. Combining with
a weighted average (weight $a_{k+1}$ for the inequality at $x^*$ and $A_k$
for the inductive hypothesis at $x_k$), and using the specific choice
$y_k = \frac{A_k x_k + a_{k+1} v_k}{A_{k+1}}$ together with the update
$v_{k+1} = v_k - \frac{a_{k+1}}{L}\nabla f(y_k)$, one obtains
after algebraic manipulations (completing the square in $v_k$):
\[
  A_{k+1}(f(x_{k+1}) - f(x^*))
  \le \frac{L}{2}\|v_0 - x^*\|^2 - \frac{L}{2}\|v_{k+1} - x^*\|^2.
\]
The algebraic details involve verifying that the momentum coefficient
$(k-1)/(k+2)$ corresponds precisely to the ratio $A_k/A_{k+1}$ that makes
$y_k$ the correct convex combination. This completes the induction.
\end{proof}

\begin{remark}[Optimality]\label{rem:optimal}
Nesterov (1983) showed that $O(1/k^2)$ is the \emph{optimal} rate for
first-order methods on the class of $L$-smooth convex functions. That is,
for any first-order method using only gradient evaluations, there exists an
$L$-smooth convex function for which the method cannot converge faster than
$\Omega(1/k^2)$. This is the celebrated information-theoretic lower bound.
\end{remark}

\begin{warning}
Nesterov acceleration does not directly apply to stochastic gradients. Naive
momentum with noisy gradients can amplify variance and diverge. This is why
SGD with momentum and Adam use different momentum schemes than classical
Nesterov acceleration.
\end{warning}

%% ============================================================
%% SECTION 7: Proximal Methods
%% ============================================================
\section{Proximal Methods}\label{sec:proximal}

\subsection{The Proximal Operator}

Many problems in ML involve non-smooth regularizers (e.g., $\ell_1$ for
sparsity). Proximal methods extend gradient descent to handle such terms.

\begin{definition}[Proximal operator]\label{def:prox}
For a closed convex function $h : \R^n \to \R \cup \{+\infty\}$, the
\emph{proximal operator} $\prox_h : \R^n \to \R^n$ is
\begin{equation}\label{eq:prox}
  \prox_h(v) = \argmin_{x \in \R^n} \left\{h(x) + \frac{1}{2}\|x - v\|^2\right\}.
\end{equation}
\end{definition}

\begin{proposition}[Well-definedness]\label{prop:prox-well}
For any closed convex $h$ with $\dom h \ne \emptyset$, the proximal operator
$\prox_h(v)$ exists and is unique for every $v \in \R^n$.
\end{proposition}

\begin{proof}
The objective $h(x) + \frac{1}{2}\|x - v\|^2$ is strongly convex (with
$\mu = 1$), hence has at most one minimizer. It is also coercive (goes to
$+\infty$ as $\|x\| \to \infty$) and lower semicontinuous, so a minimizer
exists.
\end{proof}

\begin{example}[$\ell_1$ proximal: soft thresholding]\label{ex:soft-thresh}
For $h(x) = \tau \|x\|_1$ with $\tau > 0$, the proximal operator is the
soft-thresholding operator applied componentwise:
\[
  [\prox_{\tau\|\cdot\|_1}(v)]_i = \text{sign}(v_i)\max(|v_i| - \tau, 0).
\]
\end{example}

\begin{example}[Indicator function]\label{ex:indicator-prox}
For $h = \iota_C$ (the indicator function of a closed convex set $C$, equal to
$0$ on $C$ and $+\infty$ elsewhere), $\prox_h(v) = \Pi_C(v)$, the Euclidean
projection onto $C$.
\end{example}

\subsection{Proximal Gradient Method}

Consider the composite minimization problem:
\begin{equation}\label{eq:composite}
  \min_{x \in \R^n} \; F(x) = f(x) + h(x),
\end{equation}
where $f$ is convex and $L$-smooth, and $h$ is convex but possibly non-smooth.

\begin{definition}[Proximal gradient method]\label{def:prox-grad}
The \emph{proximal gradient method} with step size $\eta = 1/L$ is:
\begin{equation}\label{eq:prox-grad}
  x_{k+1} = \prox_{\eta h}\bigl(x_k - \eta \nabla f(x_k)\bigr).
\end{equation}
\end{definition}

\begin{theorem}[Convergence of proximal gradient]\label{thm:prox-grad}
Under the composite model~\eqref{eq:composite} with $f$ convex and $L$-smooth,
the proximal gradient method with $\eta = 1/L$ achieves
\[
  F(x_k) - F(x^*) \le \frac{L\|x_0 - x^*\|^2}{2k}.
\]
The convergence rate is $O(1/k)$, matching gradient descent.
\end{theorem}

\begin{proof}
The key observation is that the proximal gradient step can be written as
$x_{k+1} = \argmin_x \{f(x_k) + \nabla f(x_k)^\top(x - x_k) + \frac{L}{2}\|x - x_k\|^2 + h(x)\}$.
By the descent lemma, $f(x_{k+1}) \le f(x_k) + \nabla f(x_k)^\top(x_{k+1} - x_k) + \frac{L}{2}\|x_{k+1} - x_k\|^2$.
Adding $h(x_{k+1})$ to both sides:
\[
  F(x_{k+1}) \le f(x_k) + \nabla f(x_k)^\top(x_{k+1} - x_k)
    + \frac{L}{2}\|x_{k+1} - x_k\|^2 + h(x_{k+1}).
\]
By the optimality condition for the proximal step, for any $x$:
\[
  f(x_k) + \nabla f(x_k)^\top(x_{k+1} - x_k) + \frac{L}{2}\|x_{k+1} - x_k\|^2 + h(x_{k+1})
  \le f(x_k) + \nabla f(x_k)^\top(x - x_k) + \frac{L}{2}\|x - x_k\|^2 + h(x).
\]
Setting $x = x^*$ and using convexity of $f$
($f(x_k) + \nabla f(x_k)^\top(x^* - x_k) \le f(x^*)$):
\[
  F(x_{k+1}) \le F(x^*) + \frac{L}{2}\|x^* - x_k\|^2.
\]
A telescoping argument identical to the one in Theorem~\ref{thm:gd-convex}
(tracking $\|x_k - x^*\|^2 - \|x_{k+1} - x^*\|^2$) yields the result.
\end{proof}

\subsection{FISTA: Fast Iterative Shrinkage-Thresholding Algorithm}

\begin{definition}[FISTA]\label{def:fista}
FISTA applies Nesterov acceleration to the proximal gradient method:
\begin{equation}\label{eq:fista}
\begin{aligned}
  y_k &= x_k + \frac{k-1}{k+2}(x_k - x_{k-1}), \\
  x_{k+1} &= \prox_{\eta h}\bigl(y_k - \eta \nabla f(y_k)\bigr),
\end{aligned}
\end{equation}
with $\eta = 1/L$ and $x_0 = x_{-1}$.
\end{definition}

\begin{theorem}[FISTA convergence]\label{thm:fista}
Under the composite model~\eqref{eq:composite}, FISTA achieves
\[
  F(x_k) - F(x^*) \le \frac{2L\|x_0 - x^*\|^2}{(k+1)^2}.
\]
This is the optimal $O(1/k^2)$ rate for first-order methods on this problem
class.
\end{theorem}

\begin{proof}
The proof follows the estimate-sequence technique of Theorem~\ref{thm:nesterov},
replacing gradient steps with proximal gradient steps. The descent lemma for the
smooth part $f$ and the convexity of $h$ allow the same inductive argument to
go through with $F = f + h$ in place of $f$.
\end{proof}

\begin{rlconnection}
Proximal methods are used in RL for imposing structure on learned policies.
For example, $\ell_1$-regularized policy optimization promotes sparse feature
representations. The mirror descent interpretation of proximal methods (with
Bregman divergences replacing the Euclidean distance) connects to natural
policy gradient, where the KL divergence plays the role of the Bregman
divergence in policy space.
\end{rlconnection}

%% ============================================================
%% SECTION 8: Interior Point Methods
%% ============================================================
\section{Interior-Point Methods}\label{sec:interior}

\subsection{Logarithmic Barrier}

Interior-point methods solve constrained problems by converting inequality
constraints into a smooth barrier term added to the objective, then following
the ``central path'' as the barrier weight decreases.

\begin{definition}[Logarithmic barrier]\label{def:barrier}
For the convex program~\eqref{eq:std-form} (with equality constraints omitted
for simplicity), the \emph{logarithmic barrier} is
\begin{equation}\label{eq:barrier}
  \phi(x) = -\sum_{i=1}^m \log(-f_i(x)),
\end{equation}
defined for $x$ in the strictly feasible set $\{x : f_i(x) < 0, \, i = 1,\dots,m\}$.
\end{definition}

\begin{proposition}[Barrier is convex]\label{prop:barrier-convex}
If $f_1, \dots, f_m$ are convex, then the logarithmic barrier $\phi$ is convex
on $\{x : f_i(x) < 0\}$.
\end{proposition}

\begin{proof}
The function $-\log(-u)$ is convex and non-decreasing on $u < 0$. Since each
$f_i$ is convex, the composition $-\log(-f_i(x))$ is convex (by the
composition rule for convex non-decreasing outer function and convex inner
function). The sum of convex functions is convex.
\end{proof}

\subsection{Central Path}

\begin{definition}[Barrier problem and central path]\label{def:central-path}
For parameter $t > 0$, the \emph{barrier problem} is
\begin{equation}\label{eq:barrier-prob}
  \min_x \; t\, f_0(x) + \phi(x).
\end{equation}
Assuming a unique minimizer $x^*(t)$ exists for each $t > 0$, the curve
$\{x^*(t) : t > 0\}$ is called the \emph{central path}.
\end{definition}

\begin{theorem}[Central path optimality bound]\label{thm:central-path}
Each point $x^*(t)$ on the central path satisfies
\begin{equation}\label{eq:central-bound}
  f_0(x^*(t)) - p^* \le \frac{m}{t},
\end{equation}
where $p^*$ is the optimal value and $m$ is the number of inequality
constraints.
\end{theorem}

\begin{proof}
At $x^*(t)$, the gradient of the barrier objective vanishes:
\[
  t\, \nabla f_0(x^*(t)) + \nabla \phi(x^*(t)) = 0.
\]
The gradient of the barrier is
$\nabla \phi(x) = \sum_{i=1}^m \frac{-\nabla f_i(x)}{-f_i(x)}
= \sum_{i=1}^m \frac{\nabla f_i(x)}{-f_i(x)}$.

Define the dual variables $\lambda^*_i(t) = \frac{1}{t(-f_i(x^*(t)))} > 0$.
Then the stationarity condition becomes
\[
  \nabla f_0(x^*(t)) + \sum_{i=1}^m \lambda^*_i(t) \nabla f_i(x^*(t)) = 0.
\]
Thus $x^*(t)$ minimizes $L(x, \lambda^*(t), 0) = f_0(x) + \sum_i \lambda^*_i(t) f_i(x)$ over $x$,
so the dual function evaluates to
\begin{align*}
  g(\lambda^*(t), 0)
  &= f_0(x^*(t)) + \sum_i \lambda^*_i(t) f_i(x^*(t)) \\
  &= f_0(x^*(t)) + \sum_i \frac{f_i(x^*(t))}{t(-f_i(x^*(t)))} \\
  &= f_0(x^*(t)) - \frac{m}{t}.
\end{align*}
By weak duality, $g(\lambda^*(t), 0) \le p^*$, so
$f_0(x^*(t)) - m/t \le p^*$, giving $f_0(x^*(t)) - p^* \le m/t$.
\end{proof}

\subsection{Barrier Method}

\begin{definition}[Barrier method]\label{def:barrier-method}
The \emph{barrier method} (or \emph{path-following method}) proceeds as follows:
\begin{enumerate}
  \item Initialize with a strictly feasible $x^{(0)}$, $t^{(0)} > 0$,
    and a growth factor $\sigma > 1$.
  \item At iteration $\ell$: solve (approximately) the barrier problem
    $\min_x \; t^{(\ell)} f_0(x) + \phi(x)$ starting from $x^{(\ell)}$,
    using Newton's method. Set $x^{(\ell+1)}$ to the result.
  \item Update $t^{(\ell+1)} = \sigma \cdot t^{(\ell)}$.
  \item Stop when $m/t^{(\ell)} \le \epsilon$.
\end{enumerate}
\end{definition}

\begin{theorem}[Complexity of the barrier method]\label{thm:barrier-complexity}
To achieve $f_0(x) - p^* \le \epsilon$, the barrier method requires
$O(\sqrt{m}\log(m/\epsilon))$ Newton steps (outer iterations), each involving
the solution of a linear system of size $n \times n$.
\end{theorem}

\begin{proof}[Proof sketch]
By Theorem~\ref{thm:central-path}, we need $t \ge m/\epsilon$, so
$\lceil \log_\sigma(m/(t^{(0)}\epsilon))\rceil$ outer iterations suffice.
The key result from the theory of self-concordant barriers (Nesterov and
Nemirovskii) is that with the optimal choice $\sigma = 1 + 1/\sqrt{m}$,
each centering step (finding $x^*(t)$ to sufficient accuracy) requires
$O(1)$ Newton steps (specifically, a constant number independent of $m$ and
$n$). The total number of outer iterations is then
$\log_\sigma(m/\epsilon) = O(\sqrt{m}\log(m/\epsilon))$,
using $\log(1 + 1/\sqrt{m}) \approx 1/\sqrt{m}$.
\end{proof}

\begin{keyresult}[title={Interior-Point Complexity}]
Interior-point methods solve convex programs with $m$ inequality constraints
to $\epsilon$-accuracy in $O(\sqrt{m}\log(m/\epsilon))$ Newton iterations.
Each Newton iteration costs $O(n^3)$ for a dense $n$-dimensional problem
(or less with structure). This polynomial complexity in $m$, $n$, and
$\log(1/\epsilon)$ makes interior-point methods highly efficient for
moderate-dimensional problems.
\end{keyresult}

\begin{rlconnection}
Interior-point methods are used to solve the linear programs that arise in the
LP formulation of finite MDPs. For large-scale problems, they complement
value iteration and policy iteration. In constrained RL (safe RL), solving
the inner constrained optimization at each policy update step can be handled
by interior-point methods when the constraint set has moderate dimension.
\end{rlconnection}

%% ============================================================
%% SECTION 9: Exercises
%% ============================================================
\section{Exercises}\label{sec:exercises}

\begin{exercise}[Dual of a linear program]\label{ex:lp-dual}
Consider the LP $\min_x \; c^\top x$ subject to $Ax = b$, $x \ge 0$.
\begin{enumerate}[label=(\alph*)]
  \item Write the Lagrangian and derive the dual problem. Show it is
    $\max_\nu \; b^\top \nu$ subject to $A^\top \nu \le c$.
  \item Prove that strong duality always holds for feasible LPs (Hint: LP
    feasible sets are polyhedral, so Slater's condition can be replaced by
    the weaker condition of feasibility).
\end{enumerate}
\end{exercise}

\begin{exercise}[KKT for least-squares with non-negativity]\label{ex:nnls}
Consider the non-negative least-squares problem:
$\min_{x \ge 0} \; \frac{1}{2}\|Ax - b\|^2$.
Write the KKT conditions explicitly and show they reduce to:
$x \ge 0$, $A^\top(Ax - b) \ge 0$, $x_i [A^\top(Ax - b)]_i = 0$ for all $i$.
\end{exercise}

\begin{exercise}[Convergence rate verification]\label{ex:gd-verify}
Let $f(x) = \frac{1}{2}x^\top Q x$ where $Q$ is symmetric positive definite
with eigenvalues $\mu = \lambda_{\min}(Q)$ and $L = \lambda_{\max}(Q)$.
\begin{enumerate}[label=(\alph*)]
  \item Verify that $f$ is $\mu$-strongly convex and $L$-smooth.
  \item Run gradient descent with $\eta = 1/L$ and verify that
    $f(x_k) \le (1 - \mu/L)^k f(x_0)$.
  \item For $Q = \diag(1, 100)$, compute the condition number and determine
    how many iterations are needed for $f(x_k)/f(x_0) \le 10^{-6}$.
\end{enumerate}
\end{exercise}

\begin{exercise}[Soft-thresholding derivation]\label{ex:soft-thresh}
Derive the soft-thresholding operator: show that
$\prox_{\tau|\cdot|}(v) = \text{sign}(v)\max(|v| - \tau, 0)$ for a scalar $v$.
(Hint: consider the three cases $v > \tau$, $|v| \le \tau$, and $v < -\tau$
separately.)
\end{exercise}

\begin{exercise}[Dual of SDP relaxation]\label{ex:sdp-dual}
Consider the boolean optimization problem $\min_{x \in \{-1,+1\}^n} x^\top W x$.
The standard SDP relaxation replaces $x x^\top$ with a matrix variable
$X \succeq 0$, $X_{ii} = 1$.
\begin{enumerate}[label=(\alph*)]
  \item Write this relaxation as a standard SDP.
  \item Derive its Lagrange dual.
  \item Show that weak duality provides a lower bound on the original
    combinatorial problem.
\end{enumerate}
\end{exercise}

\begin{exercise}[Proving the PL inequality]\label{ex:pl}
Let $f$ be $\mu$-strongly convex and differentiable.
\begin{enumerate}[label=(\alph*)]
  \item Prove that $\|\nabla f(x)\|^2 \ge 2\mu(f(x) - f(x^*))$ for all $x$.
  \item Show by example that the PL inequality can hold for non-convex
    functions (consider $f(x) = x^2 + 3\sin^2(x)$ near the origin; verify
    numerically).
  \item Prove that gradient descent converges linearly under the PL inequality
    alone (without assuming convexity).
\end{enumerate}
\end{exercise}

\begin{exercise}[FISTA for LASSO]\label{ex:fista-lasso}
The LASSO problem is $\min_x \frac{1}{2}\|Ax - b\|^2 + \tau\|x\|_1$.
\begin{enumerate}[label=(\alph*)]
  \item Identify $f(x)$ and $h(x)$ in the composite formulation.
  \item Write out the FISTA iteration explicitly (including the
    soft-thresholding step).
  \item What is the Lipschitz constant $L$ of $\nabla f$? Express it in terms
    of $A$.
\end{enumerate}
\end{exercise}

\begin{exercise}[Central path for LP]\label{ex:lp-barrier}
For the LP $\min\; c^\top x$ subject to $x \ge 0$ (i.e., $f_i(x) = -x_i$):
\begin{enumerate}[label=(\alph*)]
  \item Write the barrier function $\phi(x) = -\sum_i \log(x_i)$.
  \item Show that the central path satisfies $x^*_i(t) = 1/(tc_i)$ when $c > 0$.
  \item Verify that $c^\top x^*(t) - p^* = n/t$ where $n$ is the dimension
    and $p^* = 0$.
\end{enumerate}
\end{exercise}

\begin{exercise}[Duality gap in constrained MDP]\label{ex:cmdp}
Consider a constrained MDP with one constraint: $\max_\pi J(\pi)$ subject to
$C(\pi) \le \alpha$, where $J$ and $C$ are the expected return and expected
cost under policy $\pi$ over a finite MDP.
\begin{enumerate}[label=(\alph*)]
  \item Write the Lagrangian $L(\pi, \lambda)$ for this problem.
  \item Show that the dual function $g(\lambda) = \max_\pi L(\pi, \lambda)$
    corresponds to solving an unconstrained MDP with modified reward
    $r(s,a) - \lambda c(s,a)$.
  \item Under what condition does strong duality hold? (Hint: relate to
    Slater's condition in the occupancy-measure LP formulation.)
\end{enumerate}
\end{exercise}

\begin{exercise}[Moreau envelope and gradient]\label{ex:moreau}
The \emph{Moreau envelope} of a convex function $h$ with parameter $\gamma > 0$ is
$M_\gamma h(v) = \min_x \{h(x) + \frac{1}{2\gamma}\|x - v\|^2\}$.
\begin{enumerate}[label=(\alph*)]
  \item Show that $M_\gamma h$ is differentiable even when $h$ is not, with
    gradient $\nabla M_\gamma h(v) = \frac{1}{\gamma}(v - \prox_{\gamma h}(v))$.
  \item Prove that $M_\gamma h$ is $(1/\gamma)$-smooth.
  \item Explain how this connects to the proximal point algorithm:
    $x_{k+1} = \prox_{\gamma h}(x_k)$.
\end{enumerate}
\end{exercise}

%% ============================================================
%% SECTION 10: Summary
%% ============================================================
\section{Summary}\label{sec:summary}

\subsection{Key Theorems Recap}

The central results of this document are:
\begin{enumerate}
  \item \textbf{Weak duality} (Theorem~\ref{thm:weak-duality}): Always
    $d^* \le p^*$.
  \item \textbf{Strong duality} (Theorem~\ref{thm:strong-duality}): Under
    Slater's condition, $d^* = p^*$.
  \item \textbf{KKT conditions} (Theorems~\ref{thm:kkt-necessary}
    and~\ref{thm:kkt-sufficient}): Necessary under strong duality, sufficient
    for convex problems.
  \item \textbf{Gradient descent} (Theorems~\ref{thm:gd-convex}
    and~\ref{thm:gd-strong}): $O(1/k)$ for convex, $O(\rho^k)$ for strongly convex.
  \item \textbf{Nesterov acceleration} (Theorem~\ref{thm:nesterov}): Optimal
    $O(1/k^2)$ for smooth convex.
  \item \textbf{Interior-point methods} (Theorem~\ref{thm:barrier-complexity}):
    $O(\sqrt{m}\log(m/\epsilon))$ Newton steps.
\end{enumerate}

\subsection{Convergence Rate Summary}

\begin{table}[h]
\centering
\caption{Convergence rates for first-order and interior-point methods.}
\label{tab:rates}
\renewcommand{\arraystretch}{1.3}
\begin{tabular}{@{}llll@{}}
\toprule
\textbf{Method} & \textbf{Assumption} & \textbf{Rate} & \textbf{Iterations for $\epsilon$} \\
\midrule
Gradient descent & Convex, $L$-smooth &
  $O(L\|x_0-x^*\|^2/k)$ & $O(1/\epsilon)$ \\
Gradient descent & $\mu$-s.c., $L$-smooth &
  $O((1-\mu/L)^k)$ & $O(\kappa\log(1/\epsilon))$ \\
Nesterov accel. & Convex, $L$-smooth &
  $O(L\|x_0-x^*\|^2/k^2)$ & $O(1/\sqrt{\epsilon})$ \\
Nesterov accel. & $\mu$-s.c., $L$-smooth &
  $O((1-1/\sqrt{\kappa})^k)$ & $O(\sqrt{\kappa}\log(1/\epsilon))$ \\
Prox.\ gradient & Convex, $L$-smooth $+$ $h$ &
  $O(L\|x_0-x^*\|^2/k)$ & $O(1/\epsilon)$ \\
FISTA & Convex, $L$-smooth $+$ $h$ &
  $O(L\|x_0-x^*\|^2/k^2)$ & $O(1/\sqrt{\epsilon})$ \\
Interior point & $m$ ineq.\ constraints &
  $O(\sqrt{m}\log(m/\epsilon))$ & $O(\sqrt{m}\log(1/\epsilon))$ \\
\bottomrule
\end{tabular}
\end{table}

\begin{remark}[Lower bounds]\label{rem:lower}
The $O(1/k^2)$ rate for smooth convex functions and $O((1-1/\sqrt{\kappa})^k)$
for smooth strongly convex functions match the information-theoretic lower bounds
for first-order methods (Nesterov, 1983). These are not improvable by any
method that accesses $f$ only through gradient evaluations.
\end{remark}

\begin{remark}[Stochastic extensions]\label{rem:stochastic}
In ML and RL, we typically have access only to stochastic gradients. SGD with
step size $\eta_k = O(1/k)$ achieves $O(1/\sqrt{k})$ for convex and $O(1/k)$
for strongly convex problems (in expectation). Variance reduction techniques
(SVRG, SARAH) recover the deterministic rates up to logarithmic factors. These
stochastic methods are covered in a subsequent document.
\end{remark}

\end{document}
